\chapter{Conclusão}
\label{cap:conclusao}

Este trabalho comparou duas implementações de processamento de registros de auditoria: uma API REST que confirma o sucesso após persistir em PostgreSQL e uma API que confirma a publicação no Kafka, deixando a persistência para um consumidor. O contrato, a validação, o conteúdo determinístico e a regra de hash foram comuns. A comparação não isolou um protocolo abstrato: o broker e o consumidor adicionais fazem parte da topologia assíncrona e do uso de recursos observado.

A resposta à pergunta de pesquisa é que, no recorte experimentado, o desacoplamento temporal reduziu a latência de aceite e permitiu receber eventos durante a parada do consumidor, mas acrescentou recursos e distância entre aceite e conclusão. A síncrona utilizou menos CPU e memória e apresentou menor latência de conclusão nas condições sem falha, vinculando o produtor à disponibilidade do banco. Não se identificou uma alternativa universalmente superior. A interpretação depende de qual compromisso importa ao sistema: confirmar rapidamente o recebimento do trabalho, confirmar sua gravação ou preservar e acompanhar seu destino durante falhas.

\section{Objetivos e hipóteses}

O objetivo geral foi atendido dentro dos limites da campanha reduzida: produzir e analisar evidência experimental comparável de desempenho e resiliência. A campanha observou exatamente os 65 UUIDs planejados, em treze condições e cinco tentativas por condição, com 64 coletas completas no critério original e uma parcial qualificada. Isso não corresponde a 65 observações integralmente válidas para toda métrica. A análise de conclusão com banco parado possui quatro pares completos; os demais indicadores identificados conservam as observações elegíveis descritas no capítulo anterior.

O objetivo específico a foi atendido pela definição e verificação do modelo comum, com comparação independente de conteúdo e hash. Os objetivos b e c foram atendidos pelas variantes executáveis, seu repositório compartilhado e os fluxos distintos de confirmação. Idempotência, conflito e respostas de falha foram verificados funcionalmente; o resultado da campanha sem UUID duplicado não é utilizado como prova isolada de todos esses comportamentos. O objetivo d foi atendido pela calibração, aprovação do piloto R2, congelamento e execução dos cenários controlados, preservando as tentativas anteriores e a ordem dos pares.

O objetivo e foi atendido pela coleta conciliada de HTTP, envio--commit, vazão, erros, CPU, memória, destinos e recuperação, com a ressalva temporal da A1. A ausência de um marco não foi preenchida com horário de inserção, erro HTTP ou retomada. O objetivo f foi atendido pela discussão dos compromissos e dos resultados desfavoráveis, sem acrescentar repetições para obter um desfecho desejado. Os valores por execução, dispersões e diferenças pareadas permanecem disponíveis no conjunto público; os brutos e seus índices permitem conferir a origem das medidas.

A expectativa de menor uso básico de recursos da síncrona em carga reduzida foi corroborada descritivamente pelos pares de 12 eventos/s. A expectativa de menor latência de aceitação da assíncrona também foi corroborada nas condições comparáveis, assim como a retenção durante indisponibilidade do consumidor. A contrapartida de recursos e o maior intervalo até a conclusão apareceram nos resultados. Essa confirmação é local e descritiva, não uma demonstração estatística para uma população de aplicações.

Os aspectos financeiros e de esforço de manutenção permanecem inconclusivos: não foram medidos e não podem ser deduzidos da quantidade de componentes ou das amostras de CPU. Da mesma forma, menor tempo HTTP não demonstrou conclusão mais rápida. A indisponibilidade do PostgreSQL produziu 14 eventos terminais em DLQ na assíncrona; portanto, a interpretação de que todo aceite seria posteriormente persistido foi refutada nesse cenário. O compromisso da implementação é reter um destino de falha, não reprocessá-lo automaticamente após a volta do banco.

Recuperação também exigiu uma interpretação específica. O retorno de commits não equivale à conclusão integral do backlog, e nenhum desses marcos substitui o critério de três janelas sustentadas sob entrada ativa. Esse critério foi observado nas cinco execuções síncronas e nas cinco de consumidor parado, mas em apenas uma das cinco assíncronas de banco parado. As outras quatro não foram convertidas em sucesso pela drenagem ou pelo lag final. O indicador congelado inclui aceites destinados à DLQ na contagem sem commit; por isso, os destinos terminais também impedem seu limiar, como discutido no capítulo anterior. Não se infere falha permanente ou indisponibilidade contínua do serviço. Assim, uma superioridade geral de resiliência não foi estabelecida.

\section{Limitações e recomendações}

As conclusões se limitam a um computador Windows com WSL e Docker, cargas sintéticas, versões fixadas, um broker, uma partição e fator de replicação um. Produtor, aplicações e instrumentos compartilharam recursos físicos. Aquecimento, quotas, ordem registrada e folga dos coletores reduzem algumas fontes de interferência, mas não eliminam virtualização, caches ou processos concorrentes. Os custos da emissão de logs e da amostragem integram o ambiente observado; as médias de contêineres não representam todo o consumo da máquina.

Cinco tentativas por condição são um orçamento explicitamente reduzido, não um tamanho demonstrado suficiente por cálculo de poder. Eventos correlacionados dentro de uma coleta não ampliam esse número. A A1 foi motivada por dado já observado nas primeiras 37 tentativas e declarada antes das 28 restantes; essa decisão pós-hoc limita a conclusão, mesmo com preservação dos vereditos originais, auditoria e análise estrita de sensibilidade. A falta do instante de um commit impede conhecer sua posição na distribuição completa. A janela finita e os relógios entre processos acrescentam incerteza de observação.

Para um sistema semelhante ao protótipo, a síncrona é uma opção coerente quando a confirmação imediata de gravação e o menor caminho operacional são prioridades, desde que o produtor trate indisponibilidade e resultados incertos. A assíncrona é coerente quando se admite conclusão posterior e se deseja continuar aceitando trabalho durante parada do consumidor, desde que sejam acompanhados backlog, tentativas, DLQ e conciliação. Essas recomendações não são instruções de dimensionamento de produção e não substituem replicação no ambiente e no volume reais. Em ambas, distinguir aceite, persistência e destino terminal é mais importante do que escolher uma alternativa apenas pelo menor tempo de resposta.

\section{Trabalhos futuros}

Uma replicação pode ampliar o número de repetições e a duração da entrada ativa, com orçamento e critérios registrados antes de observar os novos dados. Máquinas separadas para gerador e serviços, novas cargas e payloads permitiriam investigar interferência e validade externa. Múltiplos brokers, replicação e falhas de rede ou de disco seriam necessários para examinar disponibilidade e durabilidade que este experimento não testou.

Outra linha é estudar políticas explícitas de reprocessamento da DLQ, limites de retry e observação de confirmação incerta, mantendo comparabilidade e sem substituir evidências desfavoráveis. Medidas específicas de energia, custo financeiro e esforço humano seriam necessárias para testar essas dimensões. A contribuição deste trabalho é, portanto, uma comparação rastreável com compromissos e limites observados, acompanhada de um protocolo e dados que permitem novas verificações, e não a eleição de uma arquitetura vencedora para qualquer contexto.
